Public opinion polling constitutes a fundamental instrument for understanding
sociopolitical demands, formulating public policies, and analyzing electoral
behavior. However, traditional field survey methods face increasing logistical
and operational challenges, including high financial costs, lengthy data
collection periods, and declining response rates. In this context, the rapid
advancement of Large Language Models (LLMs) enabled the emergence of
computational approaches directed toward human behavior simulation and synthetic
sampling (\textit{silicon sampling})~\cite{argyle_out_2023}. When conditioned on
sociodemographic attributes, such models exhibit the capacity to emulate
population attitudes and preferences at scale and on demand.

Seminal studies demonstrated that LLMs conditioned on sociodemographic profiles
can reproduce public opinion distributions with remarkable fidelity against real
empirical data~\cite{argyle_out_2023,miranda_simulating_2025}. Nevertheless, the
vast majority of existing literature focuses on North American or European
populations and on models predominantly trained in the English language. A
significant research gap remains regarding the applicability and
transferability of these methodologies to the Brazilian sociocultural context,
as well as the comparative efficacy between general-purpose global models and
models trained with an emphasis on Brazilian sociocultural and political reality.

In this paper, we investigate the ability of Brazilian Portuguese-centric
models against widely adopted global baselines to reproduce empirical public
opinion distributions in Brazil. We construct a
synthetic population comprising 2,000 personas derived from official microdata
provided by the Brazilian Institute of Geography and Statistics (IBGE),
covering demographic, economic, health, and social variables. We adopt the
survey on Brazilian perceptions of democracy conducted by the Center for Public
Opinion Research (CESOP --- \textit{Centro de Estudos de Opinião Pública da
Universidade Estadual de Campinas}) and IPEC~\cite{cesop04832} as the empirical
ground truth. We evaluate the distributional
alignment of the simulated responses using the Jensen-Shannon Divergence
(JSD)~\cite{lin_divergence_1991}.

Our primary contributions are threefold: first, we analyze an experimental
progression of four prompt engineering strategies, demonstrating that combining
feature selection, system/user role separation, Chain of Thought (CoT),
and few-shot calibration reduces JSD divergence by 77\% (from $0.525$ to $0.121$)
to achieve optimal distributional convergence; second, we conduct a systematic
architectural benchmark comparing Brazilian Portuguese-centric models
(Sabiá-4 and Sabiazinho-4) with global baselines (GPT-5-mini and LLaMA~3.2~3B),
revealing that models tailored to the Brazilian reality capture civic skepticism
with higher fidelity, whereas global models exhibit marked social desirability
and institutional optimism biases; and third, we demonstrate the utility of
Explainable Artificial Intelligence (XAI) by employing model-generated
justifications to audit persona reasoning consistency and detect cognitive
dissonances during simulation.\footnote{Replication code, prompt templates,
and evaluation pipelines are available at
\url{https://github.com/olucasaguiar/silicon-sampling-brazil}.}

The remainder of this paper is organized as follows: Section~\ref{sec:related_work}
reviews the theoretical foundation and related work; Section~\ref{sec:methodology}
details the experimental methodology, synthetic sampling procedures, and prompt
engineering scenarios; Section~\ref{sec:results} presents the quantitative
findings, bias analysis, and interpretability discussions; finally,
Section~\ref{sec:conclusion} summarizes the conclusions and outlines directions
for future research.
